\section{Schools of Magic}
Mainstream culture in Tamriel recognizes six schools of magic, each concerned with different magical effects. Beginning the game with a school of magic as a major skill will grant you starting skills from that school. All classes with magic for at least one major skill begin the game with the following spells in addition to those granted by the schools:

\begin{itemize}
	\item \textbf{Flare} (\textit{11 Magicka, Ranged, Destruction Novice}): The target takes 6 points of fire damage and an additional 1d4 on its next turn.
	\item \textbf{Heal Minor Wounds} (\textit{14 Magicka, Self, Restoration Novice}): You recover 8 health.
\end{itemize}

In combat, casting a spell requires an action unless otherwise specified.

\subsection{Alteration}
The school of Alteration covers spells that manipulate the natural world and its physical properties. Alteration magic allows you to do things like resist physical and elemental attacks, change the weight of items, open locks, and breathe underwater. Some Alteration spells target opponents with negative effects, requiring a spell attack roll.\\

If Alteration is one of your major skills, you begin the game with the following spells:
\begin{itemize}
	\item \textbf{Open Very Easy Lock} (\textit{7 Magicka, Ranged, Novice}): Automatically opens a lock of Very Easy difficulty.
	\item \textbf{Protect} (\textit{10 Magicka, Self, Novice}): +1 to your current armor rating for 5 rounds. You may cast this spell as a bonus action.
\end{itemize}

On an extreme success, Alteration attacks may have greater magnitude or duration.

\subsection{Conjuration}
\begin{wrapfigure}{L}{0.25\textwidth}
	\includegraphics[width=\textwidth]{conjurer.png}
\end{wrapfigure}

The school of Conjuration is used primarily to summon things from Oblivion. It can be used to conjure daedra and undead as well as summon daedric weapons and armor for very brief periods of time. Summoned creatures will act independently, but are allied with you. They may respond to commands if they are intelligent. You can only have one summoned minion at once. If you conjure a new minion before the current one's time runs out, the current one is dismissed. If you attack your minion, it will become hostile.\\

If Conjuration is one of your major skills, you begin the game with the following spells:
\begin{itemize}
	\item \textbf{Summon Skeleton} (\textit{45 Magicka, Self, Apprentice}): Summons a Skeleton minion for 7 rounds.
	\item \textbf{Turn Undead} (\textit{11 Magicka, Ranged, Novice}): Undead up to level 3 must use their turn to move as far from you as possible as long as they can see you, and they cannot attack you. This effect lasts until the end of the target's third turn.
\end{itemize}

\subsection{Destruction}
The school of Destruction involves all spells that directly harm the body, possessions or abilities. Destruction spells can be used to project flames, ice or lightning as well as break armor, fatigue combatants, and much more. Destruction spells almost always require a spell attack roll.\\

If Destruction is one of your major skills, you begin the game with the following spells:
\begin{itemize}
	\item \textbf{Cold Touch} (\textit{23 Magicka, Touch, Novice}): The target takes 15 points of frost damage and 7 points of stamina damage.
	\item \textbf{Shocking Touch} (\textit{14 Magicka, Touch, Novice}): The target takes 10 points of shock damage and 5 points of magicka damage.
\end{itemize}

On an extreme success, elemental damage spells get a bonus effect. Fire spells set the target on fire, dealing an additional 10\% fire damage over the course of the target's next three turns (damage is taken at the beginning of each turn). Frost spells root the target in place with ice, preventing them from moving on their next turn unless they succeed on a hard Athletics check. Shock spells jump to a nearby target within range, dealing the same effect on the new target.\\

The following are special abilities you gain access to at certain skill levels (see the mastery perk chart, section 3.2):

\begin{itemize}
\item \textbf{Dual Casting}: If you have two free hands, you can dual cast any destruction spell. Dual casted spells require double the magicka cost, but they deal triple damage.

\item \textbf{Elemental Mastery}: Choose an element: fire, frost or lightning. From now on, all spells that deal that damage type gain the corresponding effect:

\begin{itemize}
	\item \textbf{Fire}: Enemies under 20\% health are inflicted with fear for 3 turns when struck by a fire spell.
	\item \textbf{Frost}: Enemies under 20\% health are paralyzed for 3 turns when struck by a frost spell.
	\item \textbf{Lightning}: Enemies under 15\% health are immediately turned into ash when struck by a shock spell.
\end{itemize}

You may only choose one element to master, and you choose it once you achieve the mastery perk. The effects of this perk do not take effect on enemies that are resistant to that element, defined as having a resistance greater than 0.
\end{itemize}

\subsection{Illusion}
Illusion spells manipulate the senses and emotions of their targets. Illusion is useful both in and out of battle with effects such as charms, invisibility, muffling, calming, frenzy and courage. Some illusions will require spell attack rolls. People usually do not like to be tricked with illusions, so using them in social situations will require subterfuge. Illusionists must be clever!\\

If Illusion is one of your major skills, you begin the game with the following spells:
\begin{itemize}
	\item \textbf{Soothing Touch} (\textit{13 Magicka, Touch, Novice}): Targets up to level 2 are calmed for 2 rounds. Calmed targets will not move or perform any actions. They are aware of your presence, and the effect will end if you attack them.
	\item \textbf{Starlight} (\textit{14 Magicka, Self, Novice}): You emanate bright light in a 20 ft radius and dim light in a 40 ft radius for 1 minute. You may cast this spell as a bonus action.
\end{itemize}

On an extreme success, Illusion attacks may have increased magnitude or duration.

\subsection{Mysticism}
The school of Mysticism is concerned with unraveling the mysteries of the world. Spells under this school are grouped less by categorical membership characteristics and more by tradition, as these spells were considered a single school by the Psijic Order, a group of mages who created them. Effects in this school allow you to detect living creatures, dispel, reflect or absorb magical effects, harvest souls for enchanting, and move distant objects. Some Mysticism spells target opponents and will require attack rolls.\\

If Mysticism is one of your major skills, you begin the game with the following spells:
\begin{itemize}
	\item \textbf{Minor Dispel} (\textit{22 Magicka, Self, Novice}): For all magical effects currently active on you, those that were imposed by a spell that cost no more than 40 magicka are immediately ended. Applies to beneficial and harmful spells alike. Does not apply to magical effects imposed by abilities, diseases, curses, enchantments, potions, poisons or scolls.
	\item \textbf{Minor Life Detection} (\textit{15 Magicka, Self, Novice}): For the next 12 seconds, all living creatures within 60 feet of you appear to give off a faint purple glow that you can see through solid objects. Only you can see this effect.
\end{itemize}

On an extreme success, Mysticism spells may have increased magnitude or duration.

\subsection{Restoration}
Restoration magic involves all spells which heal, rest, cure, restore or fortify a target. The effects of Restoration are many and varied, and they include things such as restore health, restore stamina, fortify health, stamina and magicka, fortify skills and attributes, cure poison and diseases and resist magicka. A few Restoration spells can act as attacks, though not many.\\

\begin{figure}[h]
	\includegraphics[width=\textwidth]{healing.png}
\end{figure}

If Restoration is one of your major skills, you begin the game with the following spell:
\begin{itemize}
	\item \textbf{Absorb Health} (\textit{12 Magicka, Touch, Novice}): Transfer 5 health from the target to you.
	\item \textbf{Minor Healing Hands} (\textit{32 Magicka, Touch, Apprentice}): Heals the target for 15 health.
\end{itemize}

On an extreme success, Restoration spells may have increased magnitude or duration.